\subsection{李群的局部定义}

命题 18. 设 \(\mathbf{G}\) 是一个群，\(\mathbf{U}\) 和 \(\mathbf{V}\) 是 \(\mathbf{G}\) 中包含 \(e\) 的两个子集。假设 \(\mathbf{U}\) 具有满足下列条件的解析流形结构：

\begin{enumerate}
\def\labelenumi{(\roman{enumi})}
\item
  \(\mathrm{V} = \mathrm{V}^{-1},\mathrm{V}^2\subset \mathrm{U},\mathrm{V}\) 在 U 中是开集；
\item
  映射 \((x,y)\mapsto xy^{-1}\) 从 \(\mathbf{V}\times \mathbf{V}\) 到 \(\mathbf{U}\)，且是解析的；
\item
  对所有 \(g \in \mathbf{G}\)，存在 \(\mathbf{V}'\)，它是 \(e\) 在 \(\mathbf{V}\) 中的开邻域，使得 \(g\mathbf{V}'g^{-1} \subset \mathbf{U}\)，并且映射 \(x \mapsto gxg^{-1}\) 从 \(\mathbf{V}'\) 到 \(\mathbf{U}\)，且是解析的。
\end{enumerate}

则在 \(\mathbf{G}\) 上存在唯一一个具有下列性质的解析流形结构：

(α) 带此结构的 G 是李群；

(β) V 在 G 中是开集；

\begin{enumerate}
\def\labelenumi{(\Alph{enumi})}
\setcounter{enumi}{24}
\tightlist
\item
  G 和 U 上的流形结构在 V 上诱导出相同的结构。
\end{enumerate}

\begin{enumerate}
\def\labelenumi{(\alph{enumi})}
\item
  令 A 为 V 的开子集，\(v_{0}\) 为 V 中的元素且满足 \(v_{0}A \subset V\)。则 \(v_{0}A\) 是满足 \(v \in V\) 且 \(v_{0}^{-1}v \in A\) 的 v 所成的集合，因而是 V 的开子集（考虑 (ii)）。此外，(ii) 表明，映射 \(v \mapsto v_{0}v\) 从 A 到 \(v_{0}A\)，以及映射 \(v \mapsto v_{0}^{-1}v\) 从 \(v_{0}A\) 到 A，互为解析双射，因而是解析同构。
\item
  取 \(W\)，它是 \(e\) 在 \(V\) 中的开邻域，满足 \(W = W^{-1}\)、\(W^3 \subset V\)，并且存在图表 \((W, \phi, E)\)，它属于流形 \(U\)，其定义域为 \(W\)。对所有 \(g \in G\)，令 \(\phi_g\) 为映射 \(h \mapsto \phi(g^{-1}h)\)，它从 \(gW\) 到 \(E\)。我们证明，图表 \(\phi_g\) 是 \(G\) 的解析相容图表。令 \(g_1, g_2\) 为 \(G\) 中的元素，满足 \(g_1W \cap g_2W \neq \emptyset\)，于是 \(g_2^{-1}g_1\) 和 \(g_1^{-1}g_2\) 属于 \(W^2\)。根据 (a)，\(W \cap g_1^{-1}g_2W\) 是 \(W\) 的开子集，因此
\end{enumerate}

\[
\phi_ {g _ {1}} (g _ {1} \mathrm{W} \cap g _ {2} \mathrm{W}) = \phi (\mathrm{W} \cap g _ {1} ^ {- 1} g _ {2} \mathrm{W})
\]

是 E 的开子集 D。对于 \(d \in \mathbf{D}\)，

\[
\left(\phi_ {g _ {2}} \circ \phi_ {g _ {1}} ^ {- 1}\right) (d) = \phi \left(g _ {2} ^ {- 1} g _ {1} \phi^ {- 1} (d)\right);
\]

由 (a) 可见，\(\phi_{g_2} \circ \phi_{g_1}^{-1}\) 是解析的。

\begin{enumerate}
\def\labelenumi{(\alph{enumi})}
\setcounter{enumi}{2}
\item
  由 (b)，在 \(\mathbf{G}\) 上存在解析流形结构，使 \((\phi_g)_{g\in G}\) 成为 \(\mathbf{G}\) 上的图册。对所有 \(g_0\in \mathbf{G}\)，映射 \(g\mapsto g_0g\)（\(g\in \mathbf{G}\)）保持该图册不变，因而是带此流形结构的 \(\mathbf{G}\) 的自同构。特别地，条件 \((\mathrm{GL}_1)\) 得到满足。
\item
  令 \(v_{0} \in V\)。根据 (ii)，存在 \(A\)，它是 \(e\) 在 \(W\) 中的开邻域，且 \(v_{0}A \subset V\)。这首先证明 \(V\) 是 \(G\) 中的开集。根据 (a)，映射 \(v \mapsto v_{0}v\) 从 \(A\) 到 \(v_{0}A\) 是解析同构，其结构取自 U。根据 (c)，可见 G 和 U 上的流形结构在 \(v_{0}\mathbf{A}\) 上诱导出相同的结构，因而最终在 V 上也相同。
\item
  由 (d)、(ii) 和 (iii)，可见条件 \((\mathrm{GL}_2)\) 和 \((\mathrm{GL}_3)\) 成立。因此 G 是李群。
\item
  如果 G 上的流形结构与 G 的群结构相容，并且 V 是 G 的开子流形，那么 \((\phi_{g})_{g\in G}\) 是 G 上的图册。这就证明了命题的唯一性断言。
\end{enumerate}

命题 19. 设 G 是拓扑群，H 是李群，f 是从群 G 到群 H 的同态。假设存在 G 中 \(e_{G}\) 的一个开邻域、图表 (V, \(\phi\) , E)，它是流形 H 在 \(e_{H}\) 处的图表，以及 E 的一个具有拓扑补空间的闭向量子空间 F，使得 \(f(\mathrm{U}) \subset \mathrm{V}\)，并且 \((\phi \circ f)\mid_U\) 是从 U 到 \(\phi(\mathrm{V}) \cap \mathrm{F}\) 的同胚。则在 G 上存在唯一一个使 f 成为浸入的流形结构；该结构是 H 上流形结构在 f 下的逆像。带此结构的 G 是李群。

由于 G 的平移（相应地，H 的平移）是同胚（相应地，解析同构），f 满足《Differentiable and Analytic Manifolds》R, 5.8.1 中的条件 (R)。于是命题的前两个断言由《Differentiable and Analytic Manifolds》R, 同上处得出。考虑交换图

\[
\begin{array}{c} \mathrm{G} \times \mathrm{G} \xrightarrow {m} \mathrm{G} \\ f \times f \Big \downarrow \qquad \qquad \qquad \qquad \qquad \qquad \qquad \qquad \qquad \Big \downarrow f \\ \mathrm{H} \times \mathrm{H} \xrightarrow {n} \mathrm{H} \end{array}
\]

其中，对于 G（相应地，H）中的 x、y，\(m(x,y)=xy^{-1}\)（相应地，\(n(x,y)=xy^{-1}\)）。则 \(n\circ(f\times f)\) 是解析的，从而 \(f\circ m\) 是解析的；又因为 f 是浸入，所以 m 是解析的。因此 G 是李群。

G 上的李群结构称为 H 上李群结构在 f 下的逆像。

推论。设 G 是拓扑群，N 是 G 的离散正规子群，且 \(\pi\) 是从 G 到 G/N 的典范映射。假设 G/N 上给定了解析流形结构，并且与 G/N 的拓扑群结构相容。则在 G 上存在唯一一个使 \(\pi\) 成为浸入的流形结构；该结构是 G/N 上流形结构在 \(\pi\) 下的逆像。带此结构时，\(\pi\) 是 étale 的，G 是李群，而 G/N 是 G 按 N 作商得到的商李群。

注。设 H 是连通实或复李群，\(\hat{H}\) 是其万有覆盖\(\dagger\)，且 \(\pi\) 是从 \(\hat{H}\) 到 H 的典范映射。当我们把 \(\hat{H}\) 作为李群谈论时，始终指的是 \(\mathbf{H}\) 上的结构在 \(\pi\) 下的逆像结构。

